%!TEX root = ../thesis-1 Mids.tex
\chapter{Introduction} % (fold)
\label{cha:_introduction}

\section{Introduction} % (fold)
\label{sec1.1}

Healthcare industry is heavily reliant upon the use of networked technologies. Patients use wearable monitoring devices, implantable sensors, and remote diagnostic technologies that transmit real-time patient data to hospitals, clinicians, and cloud-based platforms. These connected technologies form what is referred to as the Internet of Medical Things (IoMT) and are central to providing contemporary patient care. Although there has been significant growth in the IoMT, there remains one major issue that the healthcare industry has yet to address: protecting the data transmitted across these systems from individuals who do not possess authorised access.

The most common attack on these systems uses stolen credentials. An attacker with a healthcare worker's login details can sign in like any other user, without breaking through any defences. Conventional access control cannot tell the legitimate owner of a credential from an impostor, because it verifies only the credential and not the person presenting it or the context of the request. After that point the system does nothing more to protect the data. The session can decrypt everything the credential holder is authorised to see, and no cryptographic control at the data layer limits what is read. In practice, a stolen login gives an attacker the same view of patient records as the clinician, including diagnoses, medication histories and mental health information.

\subsection{Attribute-Based Encryption in IoMT} % (fold)
\label{sec1.1.1}

Ciphertext-policy attribute-based encryption (CP-ABE) is among the more expressive cryptographic tools for protecting medical data. CP-ABE does not encrypt data for one named recipient. It encrypts data under a logical policy, for example Doctor AND Cardiology, and only users whose key attributes satisfy the policy can decrypt it. This fits healthcare access control, where several roles need access to overlapping data and permissions change as patients move between departments and staff rotate.

Within IoMT, the most recent and directly relevant application of CP-ABE is the scheme of Ding et al. ~\cite{ding2026badsabe}. A blockchain-enhanced framework for anonymous data sharing in 6G-enabled smart healthcare that uses distributed key generation and policy hiding. Key generation is spread across blockchain smart contracts, which removes the need for a central trusted authority, a common weak point in traditional public key infrastructures. The scheme also pairs policy hiding with anonymous authentication, so a user retrieving data reveals neither the structure of the access policy nor their own identity. Ding et al. designed the scheme for the overlapping and frequently changing roles found in modern hospitals. As it is also the most comprehensive recent work in this area, it is the primary cryptographic reference for this thesis.

The existing schemes nonetheless share a fundamental limitation with virtually all existing CP-ABE designs: the access decision is entirely static. Once a user possesses a key whose attributes satisfy the ciphertext policy, decryption is granted unconditionally. No consideration is given to whether the access attempt originates from a trusted hospital workstation or an unrecognised device in an unfamiliar location. A stolen credential carrying valid attributes will pass every verification step, indistinguishable from one used by its legitimate owner.

\subsection{Authentication Security in IoMT} % (fold)
\label{sec1.1.2}

In the Recent times the cryptographic research has sought to strengthen authentication within IoMT systems. Xu et al.\ proposed the anonymity-enhanced sequential multi-signer ring signature (ASMR) scheme, in which patients and clinicians co-sign electronic health records anonymously within separate signature rings, achieving both anonymity and traceability without dependence on a central trusted entity~\cite{xu2025asmr}. The scheme proves security in the random oracle model and achieves significant reductions in computational and communication overhead relative to prior ring signature constructions. It guarantees that only authorised parties have contributed to a given record; however, once a session is admitted, it exercises no control over what that session may subsequently decrypt.

Xu et al. proposed PQ-CLSCL, a post-quantum certificateless signcryption scheme with linkability for IoMT ~\cite{xu2025pqclscl}. Its security rests on the hardness of lattice problems, so it resists quantum adversaries. That matters in practice, because electronic health records are often kept for decades. The certificateless design avoids the key-escrow weakness of identity-based schemes and also removes the overhead of managing traditional certificates. For replay protection, the scheme binds each message to a unique session token, which lets it detect and reject replayed transmissions. This linkability check works only at the level of repeated messages. It does not consider the behavioural or environmental context of the session making the request. A stolen credential presenting a fresh, unrepeated transmission clears every verification step without raising any alarm, and PQ-CLSCL, like ASMR, exercises no influence over what an admitted session may subsequently decrypt.

\subsection{Machine Learning for IoMT Security} % (fold)
\label{sec1.1.3}

A growing body of IoMT security research uses machine learning to automate access and threat decisions. Ullah et al.\ developed a threat detection model that pairs a generative adversarial network (GAN) for handling class imbalance with a bidirectional long short-term memory network (BiLSTM) for capturing temporal attack patterns, further augmenting the system with SHAP and LIME (Local Interpretable Model-Agnostic Explanations) techniques to generate clinically interpretable explanations for each detection decision~\cite{ullah2025interpretable}. Anjum et al.\ took a complementary approach, coupling a Random Forest classifier with a blockchain-based access control layer to dynamically revise access permissions in response to observed user behaviour and access patterns~\cite{anjum2025blockchain}. Both systems deliver accurate, context-sensitive security decisions in IoMT; however, each operates exclusively at the detection or policy-decision layer with no influence on the encryption layer where data resides.

The work on cryptographic access control, authentication security and machine learning shares one gap. No existing IoMT scheme feeds a machine learning risk score into the CP-ABE access control layer, so the strength of cryptographic protection on patient data cannot change with how trustworthy an access attempt appears.

% subsection (end)

\section{Motivation}
\label{sec:motivation}

Stolen healthcare credentials give an attacker direct access to patient diagnoses, medication histories and mental health records. Data of this kind can be used for insurance fraud or extortion, or to harm a specific patient. No mechanism currently deployed in IoMT stops this exposure at the cryptographic layer. Anyone holding a healthcare worker's credentials gets the same decryption rights as the worker, however unlikely the circumstances of the request. A login from another country, on a device the system has never seen, at an unusual hour, still meets no barrier at the data layer. Intrusion detection systems can flag behaviour like this, but they have no way to withhold the data. Cryptographic controls do protect the data, yet they take no account of how a credential is being used.

Embedding a login risk score as an attribute within a CP-ABE session key makes the encryption layer risk-aware. A session assessed as low-risk receives a key capable of decrypting sensitive medical records; one assessed as high-risk (for instance, because the login originated from an unfamiliar country, device, and network simultaneously) receives a restricted key that cannot access sensitive data, regardless of the credential's cryptographic validity. The protection sits on the encrypted data, not at the login prompt. An attacker who has stolen a second factor still cannot get past it without also reproducing the legitimate user's login context. This thesis designs and implements the system and then evaluates it.

\section{Problem Statement}
\label{sec:problemstatement}

Current IoMT access control schemes give decryption rights to any session that presents valid credential attributes. They treat a routine request from a known device the same way as a login from an unfamiliar country or network on a device never seen before. Authentication behaves similarly. Once a session is admitted, the system discards its contextual information and treats the credential as fully legitimate. A stolen credential therefore opens the same data as one used by its rightful owner. To the best of the author's knowledge, no existing IoMT mechanism addresses this risk at the cryptographic layer, where patient data is actually protected. Nothing stands between the theft of a credential and full access to the data it unlocks.

\section{Research Questions}
\label{sec:researchquestions}

This thesis addresses three research questions:

\begin{enumerate}
    \item How well can a machine learning model trained on login context and user history features separate credential misuse from legitimate logins, compared with a statistical risk-based authentication baseline?
    \item How effectively does embedding a machine learning risk level as an attribute in CP-ABE session keys prevent credential misuse from decrypting sensitive IoMT records, and what false challenge rate does it impose on legitimate users?
    \item What computational overhead does risk-adaptive CP-ABE add compared with static CP-ABE?
\end{enumerate}

\section{Objectives}
\label{sec:objectives}

The objectives of this research, each corresponding directly to one research question, are as follows:

\begin{enumerate}
    \item To develop and evaluate machine learning models that detect credential-misuse login attempts from login-context and user-history features, and to compare their performance with a statistical risk-based authentication baseline.
    \item To design and implement a risk-adaptive CP-ABE scheme in which short-lived session keys carry a three-level risk attribute, and patient records are protected by access policies tiered by data sensitivity. The scheme's attack block rate and false challenge rate will be measured under simulated credential misuse
    \item To compare the computational overhead of the risk-adaptive scheme with a static CP-ABE baseline, measuring key generation, encryption, decryption and session-key reissue times.
\end{enumerate}

\section{Research Methodology}
\label{sec:researchmethodology}

The research has three phases, one for each component of the proposed system. The first phase constructs and evaluates the machine learning risk model; the second integrates the trained risk model with the CP-ABE access control layer; and the third evaluates the integrated system under controlled experimental conditions. Each phase produces an independently testable artefact that serves as the input to the next. The first phase addresses the first research question, the second and third phases together address the second, and the cryptographic benchmarks of the third phase address the third.

In the first phase of this research is dedicated to developing the risk model. A user-level random sample is drawn from the publicly available RBA dataset of Wiefling et al.~\cite{wiefling2023pump}, in which all login attempts of the selected users are retained so that complete login histories are preserved. The dataset contains over 33 million login attempts synthesised from the real login behaviour of 3.3 million users of a large online service. To ensure temporal validity, data are partitioned by user and time, so that models are evaluated exclusively on later login attempts from accounts withheld during training. The feature set combines login context with user history. The login context signals are the country and region derived from the IP address, autonomous system number, browser, operating system, device type, round-trip time and login hour. The user history features record whether the country, device or network is being seen for the first time on that account. The risk model is evaluated on the attack labels supplied with the dataset. It is not evaluated on the simulated attacks from the third phase, because those attacks are built from the model's own features and would inflate its scores. LightGBM is the main risk-scoring model. XGBoost provides a second supervised baseline, and Isolation Forest covers unsupervised anomaly detection. Cost-sensitive learning handles the class imbalance, and every performance metric is reported with a 95% bootstrap confidence interval.

The second phase connects the trained risk model to the CP-ABE access control layer. Risk scores are discretised into three attribute levels (low, medium and high), and these levels are embedded in session-level decryption keys built with the Charm-Crypto framework. Session keys are short-lived and are reissued when a session's risk level changes, because CP-ABE keys cannot be efficiently revoked once issued; the cost of this reissue is measured in the third phase. A standard hybrid encryption design is adopted: AES-GCM encrypts each medical record directly, while CP-ABE is applied only to the AES key, limiting the cryptographic overhead. Synthetic patient records produced using Synthea are organised into three sensitivity tiers, each of which requires a different minimum risk attribute for successful decryption.

In third phase of this research evaluates the integrated system under controlled experimental conditions. Credential misuse is simulated by replaying authenticated account histories from login contexts (countries, devices, and networks) not previously associated with those accounts. To avoid an evaluation that favours the model by construction, the simulation also includes harder misuse scenarios in which the attacker reproduces part of the legitimate user's context, such as the same country or a previously used browser, and the attack block rate is reported separately for each scenario. System performance is assessed on four metrics: the attack block rate, the false challenge rate imposed on legitimate users, session-key reissue time, and the cryptographic overhead relative to a static CP-ABE baseline.

\section{Thesis Organisation}
\label{sec:section_ThesisOrg}

The rest of the thesis is organised as follows. Chapter 2 reviews the literature on cryptographic access control and authentication schemes for IoMT, along with machine learning approaches to security in this domain, and ends by identifying the research gap this thesis addresses. Chapter 3 describes the proposed methodology, covering dataset preparation, feature engineering, model training and evaluation, CP-ABE policy design, and the integration prototype. Chapter~4 reports the experimental results, including the risk model evaluation, the credential misuse simulation, and the cryptographic performance benchmarks. Chapter~5 concludes the thesis with a summary of the principal findings, a discussion of limitations, and recommendations for future research directions.

\section{Chapter Summary}
\label{sec:section_Sum}

This chapter has introduced the problem of credential misuse in IoMT, wherein stolen credentials can be used to decrypt sensitive patient records because existing CP-ABE access control schemes are static in nature and current authentication mechanisms exercise no influence over the decryption layer. This chapter explained why machine learning risk scoring should be built into CP-ABE and set out the problem formally. It then gave the research questions and objectives, described the three-phase methodology, and outlined the structure of the thesis. The following chapter reviews the body of literature that informs and contextualises this work.
